\chapter{Xi Scale-Defect Antisymmetry on the Centered Uroboros Cell}
\label{chap:xi-scale-defect}

\status{PROVED as a consequence of the completed xi functional equation and the explicit scale quotient. No zero-location or RH claim is made.}

Chapter~\ref{chap:uroboros-torus} introduced the explicit dilation quotient associated with the scale \(32=2^5\).  The present chapter asks whether the transformed completed xi function itself can be periodic under that dilation.  The answer is no; the exact Riemann symmetry instead produces an antisymmetric scale defect and a reciprocal fundamental cell centered on the self-dual point.

\section{Dilation and inversion}

Let
\[
X(u)=\xi\!\left(\frac{u}{1+u}\right),\qquad u>0.
\]
The xi functional equation gives
\[
\boxed{X(u)=X(1/u).}
\]

For a fixed scale \(a>1\), define
\[
T_a(u)=au,\qquad J(u)=\frac1u.
\]
Then
\[
JT_aJ(u)=\frac{u}{a}=T_a^{-1}(u),
\]
so
\[
\boxed{JT_aJ=T_a^{-1}.}
\]
Thus inversion normalizes the dilation subgroup \(a^{\mathbb Z}\) and therefore descends to the quotient
\[
\mathbb C^*/a^{\mathbb Z}.
\]

For \(a=32\), a reciprocal radial fundamental cell is
\[
\boxed{32^{-1/2}\le |u|\le32^{1/2}.}
\]
The ordinary Riemann inversion \(u\mapsto1/u\) exchanges the two boundaries and fixes the central positive point \(u=1\), corresponding to \(s=1/2\).

\section{Exact scale-defect antisymmetry}

Define
\[
\Delta_a(u)=X(au)-X(u).
\]
Then
\[
\begin{aligned}
\Delta_a\!\left(\frac1{au}\right)
&=X(1/u)-X(1/(au))\\
&=X(u)-X(au)\\
&=-\Delta_a(u).
\end{aligned}
\]
Hence:

\begin{theorem}[SOH-G008 scale-defect antisymmetry]
For every \(a>1\) and every positive \(u\),
\[
\boxed{\Delta_a(1/(au))=-\Delta_a(u).}
\]
The defect involution \(u\mapsto1/(au)\) has the unique positive fixed point \(u=a^{-1/2}\), so
\[
\boxed{\Delta_a(a^{-1/2})=0}
\]
and therefore
\[
\boxed{X(a^{1/2})=X(a^{-1/2}).}
\]
\end{theorem}

For the Uroboros scale,
\[
\boxed{X(\sqrt{32})=X(1/\sqrt{32}).}
\]
This is not a scale-periodicity assumption.  It follows from the ordinary reciprocal symmetry because the endpoints are reciprocal.

\section{Logarithmic centering}

Write
\[
\lambda=\log u,\qquad Y(\lambda)=X(e^\lambda),\qquad L=\log a.
\]
Then
\[
Y(-\lambda)=Y(\lambda),
\]
while
\[
d_L(\lambda)=Y(\lambda+L)-Y(\lambda)
\]
satisfies
\[
\boxed{d_L(-L-\lambda)=-d_L(\lambda).}
\]
The reciprocal fundamental interval
\[
[-L/2,L/2]
\]
is centered at \(\lambda=0\), exactly the self-dual point \(u=1\leftrightarrow s=1/2\).

\section{No-go for exact scale periodicity}

Suppose that for some \(a>1\),
\[
X(au)=X(u)
\]
held for every positive \(u\).  Iteration would give
\[
X(u)=X(a^n u).
\]
As \(n\to\infty\), the corresponding Riemann coordinate approaches \(s=1\), and continuity of \(\xi\) gives
\[
X(u)=\xi(1)=\frac12
\]
for every positive \(u\).  Hence \(\xi\) would be constant on the interval \((0,1)\) and therefore constant everywhere by the identity theorem, contradicting the nonconstancy of the completed xi function.  Consequently
\[
\boxed{X(au)=X(u)\text{ globally is impossible for every }a>1.}
\]

Likewise, if
\[
X(au)=cX(u)
\]
with a constant multiplier \(c\), then taking \(u\to\infty\) gives \(1/2=c/2\), hence \(c=1\), reducing to the impossible periodic case.

\section{What survives of the torus picture}

The completed xi function does not descend to the Uroboros quotient as an ordinary scale-periodic function.  What survives exactly is:
\begin{itemize}
\item the quotient itself, once the scale identification is imposed;
\item the descended reciprocal involution because \(JT_aJ=T_a^{-1}\);
\item the self-dual centered cell;
\item the antisymmetric scale defect \(\Delta_a\);
\item the exact equality of xi values at reciprocal cell boundaries.
\end{itemize}

\section{Proof firewall}

No nonconstant factor of automorphy
\[
X(au)=M_a(u)X(u)
\]
has been independently derived.  Defining \(M_a\) as the quotient \(X(au)/X(u)\) would be tautological away from zeros and is not promoted as a proof mechanism.  No zero-location result follows from SOH-G008 alone.  SOH-G003 and the Riemann Hypothesis remain OPEN.
